\documentclass[10pt]{article}

\usepackage[utf8]{inputenc}

\usepackage[T1]{fontenc}

\usepackage{amsmath}

\usepackage{amsfonts}

\usepackage{amssymb}

\usepackage[version=4]{mhchem}

\usepackage{stmaryrd}

\usepackage{mathrsfs}

\usepackage{graphicx}

\usepackage[export]{adjustbox}

\graphicspath{ {./images/} }



\begin{document}

мерение квантовой наблюдаемой $X$ с чисто точечным спектром, то состояние системы меняется по правилу



\[

\rho \rightarrow E(x) \rho E(x) / \operatorname{Tr} \rho E(x)

\]



с вероятностью $\operatorname{Tr} \rho E(x)$, где $\rho$ - оператор плотности состояния перед измерением, $x$ - измеренное значение, $E(x)-$ соответствующий спектральный проектор наблюдаемой $X$.



Попытки сформулировать постулат редукции для наблюдаемых с непрерывным спектром встречает принципиальные трудности (см. [2]-[4]). Для описания изменения состояния при произвольном квантовом измерении вводится понятие инструмента, охватывающее неточные измерения, не удовлетворяющие В. Г. (см. [2]). При этом В. г. и постулат редукции рассматриваются как дополнительные требования, выражающие свойства точности и мйимальной возмущаемости и не входящие в исходные посылки теории квантового измерения. Инструмент $\{\mathscr{\mu}(B) ; B \in \mathscr{B}(\mathscr{X})\}$ с множеством исходов $\mathscr{X}$ удовлетворяет условию вос п р о и з водимости, если



\[

\operatorname{Tr} \mathscr{\mu}(B)[\mu(C)[\rho]]=\operatorname{Tr} \cdot \mu(B \cap C)[\rho], B, C \in \mathscr{B}(\mathscr{X}),

\]



для любого оператора плотности р. Доказывается, что вполне положительный инструмент, удовлетворяющий этому условию, с необходимостью имеет дискретное множество возможных исходов (см. [3]). Условие воспроизводимости для измерений наблюдаемых с непрерывным спектром может быть удовлетворено лишь путем включения в понятие квантового измерения предельных процедур, в результате к-рых система переходит в нек-рое алгебраич. состояние, не являющееся нормальным и не задаваемое оператором плотности (см. [4]).



Лum.: [1] Н е й ман фон Д ж., Математические основы квантовой механики, пер. с нем., М., 1964; [2] Davies E. B., Quantum theory of open systems, L.- N.Y., 1976; [3] O zawa M., «J. Math. Phys.», 1984, v. 25, p. 79-87; [4] S ri nivas M. D., "Comm. Math. Phys.», 1980, v. 71, p. 131-58. A. $C$. Холево.



ВПОЛНЕ ПОЛОХКИЕЛЬНОЕ ОТОБРАХЕНИЕ - ОТОбражение $\Phi$ из $C^{*}$-алгебры $\mathscr{A}$ в $C^{*}$-алгебру $\mathscr{B}$, удовлетворяющее условию: для любого $n=1,2, \ldots$ и любой положительной $(n \times n)$-матрицы с элементами $X_{i j} \in \mathscr{\varkappa}, i, j=1, \ldots, n$, матрица с элементами $\Phi\left(X_{i j}\right) \in \mathscr{B}$ положительна. Всякое В. п. о. положительно, то есть $X \geqslant 0$ влечет $\Phi(X) \geqslant 0$; пример положительного отображения, не являющегося В. П.о.,- транспонирование в алгебре комплексных $(d \times d)$-матриц $(d>1)$. Всякий *-гомоморфизм и условное ожидание являются В. П.о. Если или я коммутативна, то всякое линейное положительное отображение из $\mathscr{A}$ в $\mathscr{B}$ является В. п. о.



Структура В. П.о. раскрывается те о ремой Стайн с п ри нга: пусть $\Phi-$ линейное В. П. о. из $C^{*}$-алгебры $\mathscr{A}$ в алгебру $\mathfrak{B}(H)$ всех ограниченных операторов в гильбертовом пространстве $H$, тогда существует гильбертово пространство $K \supset H,{ }^{*}$-гомоморфизм $\pi$ алгебры $\mathscr{A}$ в $\mathfrak{B}(K)$ и ограниченный оператор $V$ из $H$ в $K$ такие, что $\Phi(X)=V^{*} \pi(X) V$. В частности, если $\mathscr{A}=\mathfrak{B}(H)$, то всякое нормальное линейное В. П. о. имеет вид:



\[

\Phi(X)=\sum_{i} V_{i}^{*} X V_{i}, \quad X \in \mathfrak{B}(H),

\]



где $V_{i} \in \mathfrak{B}(H)$ и ряд $\sum_{i} V_{i}^{*} V_{i}$ сходится в сильной операторной топологии $\mathfrak{B}(H)$.



Понятие В. п.о. играет важную роль в структурной теории алгебр Неймана, в алгебраич. формулировке неравновесной квантовой статистич. механики и в теории квантового измерения. См. Динамическая полугруппа, Инструмент.



Jum.: [1] Arveron W., "Acta math.», 1969, v. 123, p. 141-224; [2] Stinespring W. F., "Proc. Amer. Math. Soc.», 1955, v. 6, p. 211-16; [3] Størmer E., «Lect. Notes in Phys.», 1974, v. 29, $\begin{array}{ll}\text { p. 85-106. } & \text { А. Х. Холево. }\end{array}$



ВПОЛНЕ ПРИВОДИМОЕ ПРЕДСТАВЛЕНИЕ Г Р У П П $G$ в вектор ном $\Pi$ р ост ра нс т в $V-$ представление $T$, любое инвариантное подпространство которого обладает ортогонатьным дополнением, то есть для любого инвариантного относительно $T$ подпространства $V_{1}$ в $V$ существует инвариантное подпространство $V_{2}$ в $V$ такое, что $V$ есть прямая сумма $V_{1}$ и $V_{2}$. Представления конечных групп, компактных групп, а также унитарные представления вполне приводимы. Ю. А. Неретин.



ВРАШАТЕЛЬНОЕ ДВИХЕНИЕ т в Р Д О го те Л а.



\begin{enumerate}

  \item В. Д. вокруг не под вижной оси - движение твердого тела, при котором все его точки, двигаясь в параллельных плоскостях, описывают окружности с центрами, лежащими на одной неподвижной прямой, называемой о с ь в в ше н и я. Тело, совершающее В. Д., имеет одну степень свободы, и его положение относительно данной системы отсчета определяется углом поворота $\varphi$ между неподвижной полуплоскостью и полуплоскостью, жестко связанной с телом, проведенными через ось вращения (рис. 1).

\end{enumerate}



Pис. 1. Врашательное движение тела вокруг неподвижной оси.



\begin{center}

\includegraphics[max width=\textwidth]{2023_07_08_a185f34fbfdb2aad6b4ag-1}

\end{center}



В. д. задается уравнением $\varphi=f(t)$, где $t-$ время. Основные кинематич. характеристики В. Д. тела: его уаловая скорость



\[

\omega=d \varphi / d t

\]



и угловое ускорение



\[

\varepsilon=d \omega / d t=d^{2} \varphi / d t^{2}

\]



Для любой точки тела, находящейся на расстоянии $h$ от оси вращения, л и н ей н а я с кор ост ь



\[

v=\omega h

\]



касательное ускорение



\[

w_{\tau}=h \varepsilon

\]



нормальное ускорение



\[

w_{n}=h \omega^{2}

\]



полное ускорение



\[

w=h \sqrt{\varepsilon^{2}+\omega^{4}}

\]



Таким образом, скорости и ускорения всех точек тела пропорциональны их расстояниям от оси вращения.



Основными динамич. характеристиками В. д. тела являются его главные моменты количеств движения относительно связанных с телом осей $x, y, z$ ( $z-$ ось вращения), равные:



\[

K_{x}=-I_{x z} \omega, K_{y}=-I_{y z} \omega, K_{z}=I_{z} \omega

\]



и кинетич. энергия



\[

T=1 / 2 I_{z} \omega^{2},

\]



где $I_{z}$ - осевой, а $I_{x z}, I_{y z}$ - центробежные моменты инерции.



\begin{enumerate}

  \setcounter{enumi}{1}

  \item В. Д. вокруг точки (или сферическое дви же ние) - движение твердого тела, имеющего одну неподвижную точку $O$ (напр., Движение гироскопа, закрепленного в кардановом подвесе). Каждая из точек тела при этом В. д. перемещается по поверхности сферы с центром в точке $O$. В. Д. тела вокруг точки слагается из серии элементарных или мгновенных В. Д. вокруг мгновенных осей вращения, проходящих через эту точку. Мгновенная ось вращения непрерывно изменяет свое положение как по отношению к системе отсчета, в к-рой рассматривается движение тела, так и в самом теле, образуя при этом две конич. поверхности, называемые соответственно не под вижным и под в иж ны м а к с о да м и. Качением подвижного аксоида

\end{enumerate}



\end{document}